\documentclass[11pt]{article}
\usepackage[letterpaper,margin=1in]{geometry}
\usepackage{amsmath,amssymb,amsthm,mathtools}
\usepackage{graphicx}
\usepackage{microtype}
\usepackage{xcolor}
\usepackage{enumitem}
\usepackage[hidelinks]{hyperref}

\newtheorem{theorem}{Theorem}
\newtheorem{lemma}{Lemma}
\newcommand{\R}{\mathbb R}
\newcommand{\N}{\mathbb N}
\newcommand{\erfc}{\operatorname{erfc}}
\newcommand{\interior}{\operatorname{int}}
\newcommand{\cH}{\mathcal H}

\title{Infinite ultrathin annuli refute the universal\\
complementary-error-function plunge profile}
\author{Automated mathematics run \texttt{fa\_banach\_001}\\
Lane 10 --- candidate counterexample for expert review}
\date{9 August 2026}

\begin{document}
\maketitle

\begin{center}
\fbox{\parbox{0.91\textwidth}{
\textbf{Status.} Candidate full counterexample to Conjecture 2.6 of
Halvdansson.  There is a compact regular-closed radial set with finite boundary
length for which the asserted uniform $O(R^{-1})$ eigenvalue profile fails.
The proof is complete subject to verification of the source's standard
uniform-in-$k$ interpretation of the displayed profile.
}}
\end{center}

\section{Source conjecture}

For the standard Gaussian window
$g_0(t)=2^{1/4}e^{-\pi t^2}$, let $A_{\Omega}^{g_0}$ be the time-frequency
localization operator with symbol $\Omega\subset\R^2$, and write
$\lambda_k^{\Omega}$ for its decreasingly ordered eigenvalues.  Conjecture
2.6 of Halvdansson \cite{Halvdansson} states that, for every compact,
regular-closed $\Omega$ of finite boundary length,
\begin{equation}\label{eq:source-conj}
 \sup_k\left|\lambda_k^{R\Omega}
 -\frac12\erfc\!\left(
 \sqrt{2\pi}\frac{k-R^2|\Omega|}{R|\partial\Omega|}
 \right)\right|=O(R^{-1}).
\end{equation}
The supremum records the uniform-in-$k$ reading used throughout the source's
plunge-profile results and numerical $L^\infty$ comparisons.  Without this
uniformity, the formula does not describe the moving plunge region
$k-R^2|\Omega|=O(R)$.

\begin{figure}[ht]
  \centering
  \includegraphics[width=0.97\textwidth]{figures/open_problem_crop.png}
  \caption{Conjecture 2.6 and its eigenvalue-counting interpretation on
  printed page 13 of the source paper.}
\end{figure}

\section{Proof intuition}

For a radial symbol, all localization operators are diagonal in the Hermite
basis.  After the change of variables $s=\pi|z|^2$, the eigenvalue attached to
the $k$th Hermite function is exactly the integral of the Gamma density
$e^{-s}s^k/k!$ over the radial symbol.  Thus a radial annulus is a short
interval sampled by a probability density of width comparable to the square
root of its location.

We use infinitely many concentric annuli.  Their radii decrease geometrically,
so the sum of the lengths of all boundary circles is finite, but their radial
widths decrease inductively much faster.  At a selected scale $R_m$, the
$m$th annulus has scaled boundary length tending to infinity, whereas its
entire Gamma-density bump has height $o(1)$.  Later annuli have negligible
total symbol area at that scale.  Earlier annuli form a finite radial union,
so the source's proved finite-component profile applies to them.

At any fixed threshold $0<\theta<1/2$, the full operator therefore has almost
the same number of eigenvalues above $\theta$ as the finite prefix.  The
conjectured formula instead charges the invisible tail for all of its boundary
circles, creating an unbounded discrepancy.  A uniform $O(R^{-1})$
eigenvalue error could create only an $O(1)$ counting discrepancy, which is
the contradiction.

\section{Radial diagonalization and two lemmas}

Put
\[
 p_k(s):=e^{-s}\frac{s^k}{k!}\qquad(k=0,1,\ldots).
\]
If $E\subset[0,\infty)$ and
\[
 \Omega_E:=\{z\in\R^2:\pi|z|^2\in E\},
\]
then the Hermite functions diagonalize $A_{R\Omega_E}^{g_0}$, and its
unordered eigenvalues are
\begin{equation}\label{eq:gamma-eigs}
 \mu_k(R,E)=\int_{R^2E}p_k(s)\,ds.
\end{equation}
Indeed,
$|V_{g_0}h_k(z)|^2=e^{-\pi|z|^2}(\pi|z|^2)^k/k!$, and
$s=\pi|z|^2$ transforms planar radial measure into $ds$.  This is the same
radial formula used in the disk and annulus calculations of the source
\cite[Section 2]{Halvdansson}.

\begin{lemma}[A thin annulus has a small spectral bump]\label{lem:bump}
There is an absolute constant $C_0$ such that, whenever $a>0$, $\ell>0$,
$R^2a\geq1$, and $I=[a,a+\ell]$,
\[
 \sup_{k\geq0}\int_{R^2I}p_k(s)\,ds
 \leq C_0\frac{R\ell}{\sqrt a}.
\]
\end{lemma}

\begin{proof}
For fixed $s$, $p_k(s)$ is the probability that a Poisson random variable of
mean $s$ equals $k$.  Unimodality in $k$ and Stirling's lower bound for
$\lfloor s\rfloor!$ give the standard maximal-atom estimate
\[
 \sup_{k\geq0}p_k(s)\leq C_0s^{-1/2}\qquad(s\geq1).
\]
The interval $R^2I$ has length $R^2\ell$ and lies to the right of $R^2a$.
Integrating the displayed estimate gives
$C_0R^2\ell/(R\sqrt a)$.
\end{proof}

For $A,P>0$, define the model profile
\[
 F_{R,A,P}(x):=\frac12\erfc\!\left(
 \sqrt{2\pi}\frac{x-R^2A}{RP}\right),
 \qquad
 b_\theta:=\frac{\erfc^{-1}(2\theta)}{\sqrt{2\pi}}.
\]
Thus $b_\theta>0$ for $0<\theta<1/2$.

\begin{lemma}[Uniform profile implies bounded counting error]
\label{lem:counting}
Let $(\lambda_k(R))$ be decreasing.  If
\[
 \sup_k|\lambda_k(R)-F_{R,A,P}(k)|\leq C/R,
\]
then, for each fixed $\theta\in(0,1)$,
\begin{equation}\label{eq:counting-law}
 \#\{k:\lambda_k(R)>\theta\}
 =R^2A+RPb_\theta+O(1).
\end{equation}
The $O(1)$ constant may depend on $C,P,$ and $\theta$, but not on $R$.
\end{lemma}

\begin{proof}
The count is sandwiched between the numbers of integers $k$ for which
$F_{R,A,P}(k)>\theta+C/R$ and
$F_{R,A,P}(k)>\theta-C/R$.  Solving these two inequalities for $k$ gives
\[
 k<R^2A+RP\frac{\erfc^{-1}(2\theta+O(R^{-1}))}{\sqrt{2\pi}}.
\]
The inverse complementary error function is continuously differentiable near
$2\theta$.  Its perturbation is $O(R^{-1})$, which changes the cutoff by
$O(1)$.  Integer rounding contributes another $O(1)$.
\end{proof}

\section{The infinite-annulus symbol}

Set
\[
 a_m:=4^{-m},\qquad c_m:=a_m\qquad(m\geq1).
\]
We choose positive lengths $\ell_m<a_m/10$ inductively.  Put
\begin{equation}\label{eq:Rm}
 R_m:=\frac{c_m\sqrt{a_m}}{\ell_m}.
\end{equation}
The lengths can be chosen to have all of the following properties:
\begin{align}
 R_m&\longrightarrow\infty,
 &R_m\sqrt{a_m}&\longrightarrow\infty, \label{eq:growth}\\
 R_i^2\sum_{j>i}\ell_j&\leq c_i
 &&(i\geq1), \label{eq:tail-area}
\end{align}
and the finite-prefix counting constants introduced below are
$o(R_m\sqrt{a_m})$.

Here is the elementary induction.  After $\ell_1,\ldots,\ell_{m-1}$ have
been chosen, the union of the first $m-1$ annuli is fixed.  Proposition 2.5
of the source gives its uniform $O(R^{-1})$ profile, and Lemma
\ref{lem:counting} gives a finite counting constant, uniformly for thresholds
in a compact subinterval of $(0,1)$.  Choose $\ell_m$ small enough that the
two quantities in \eqref{eq:growth} exceed any prescribed bounds and that this
fixed finite-prefix constant is at most
$2^{-m}R_m\sqrt{a_m}$.  Also impose
\[
 \ell_m\leq2^{-m}\min_{i<m}\frac{c_i}{R_i^2}.
\]
For each fixed $i$, summing this last condition over future $m$ yields
\eqref{eq:tail-area}.  All requirements are compatible because $\ell_m$ may
be arbitrarily small.

Define
\begin{equation}\label{eq:E-Omega}
 E:=\{0\}\cup\bigcup_{m\geq1}[a_m,a_m+\ell_m],
 \qquad \Omega:=\Omega_E.
\end{equation}
The intervals are disjoint and accumulate only at $0$.  Hence $E$ and
$\Omega$ are compact.  The interior of $\Omega$ is the union of the open
annuli, and its closure is exactly $\Omega$; thus $\Omega$ is regular closed.
Its area and boundary length are
\begin{equation}\label{eq:geometry}
 |\Omega|=\sum_m\ell_m<\infty,
 \qquad
 |\partial\Omega|
 =2\sqrt\pi\sum_m\bigl(\sqrt{a_m}+\sqrt{a_m+\ell_m}\bigr)<\infty.
\end{equation}
The accumulation point at the origin has zero $\cH^1$ measure.  In
particular, \eqref{eq:E-Omega} satisfies every geometric hypothesis stated in
Conjecture 2.6.

\section{Failure of the conjectured counting law}

Fix $\theta\in(0,1/2)$, for example $\theta=1/4$.  Let
$E_{<m}:=\bigcup_{j<m}[a_j,a_j+\ell_j]$, and write $A_{<m}$ and $P_{<m}$
for the area and boundary length of $\Omega_{E_{<m}}$.  This is a finite
union of annuli.  Proposition 2.5 of the source and Lemma
\ref{lem:counting} give, at $R=R_m$ and uniformly for
$u\in[\theta/2,\theta]$,
\begin{equation}\label{eq:prefix-count}
 N_{<m}(u)
 =R_m^2A_{<m}+R_mP_{<m}b_u+O_m(1),
\end{equation}
where the construction makes $O_m(1)=o(R_m\sqrt{a_m})$.

Let $\tau_{m,k}$ be the contribution in \eqref{eq:gamma-eigs} of the annuli
with indices $j\geq m$.  Since
$R_m\ell_m/\sqrt{a_m}=c_m$ and $R_m^2a_m\geq1$ for large $m$, Lemma
\ref{lem:bump} bounds the $m$th contribution by $C_0c_m$.  For the later
annuli we use only $p_k\leq1$ and \eqref{eq:tail-area}:
\[
 \sup_k\sum_{j>m}\int_{R_m^2[a_j,a_j+\ell_j]}p_k(s)\,ds
 \leq R_m^2\sum_{j>m}\ell_j\leq c_m.
\]
Consequently
\begin{equation}\label{eq:tail-norm}
 0\leq\tau_{m,k}\leq\delta_m:=(C_0+1)c_m
 \qquad\text{for every }k.
\end{equation}

Because all pieces are diagonal in the same Hermite basis,
\eqref{eq:tail-norm} immediately gives
\begin{equation}\label{eq:sandwich}
 N_{<m}(\theta)\leq N_{\Omega,R_m}(\theta)
 \leq N_{<m}(\theta-\delta_m)
\end{equation}
for all large $m$.  Since $b_u$ is Lipschitz near $\theta$,
\eqref{eq:prefix-count} and \eqref{eq:sandwich} imply
\begin{equation}\label{eq:actual-upper}
 N_{\Omega,R_m}(\theta)
 \leq R_m^2A_{<m}+R_mP_{<m}b_\theta
 +C_1R_mP_{<m}c_m+o(R_m\sqrt{a_m}).
\end{equation}

On the other hand, the count forced by the conjectured profile is
\[
 Q_m:=R_m^2|\Omega|+R_m|\partial\Omega|b_\theta.
\]
The $m$th annulus alone has boundary length
\[
 2\sqrt\pi\bigl(\sqrt{a_m}+\sqrt{a_m+\ell_m}\bigr)
 \geq4\sqrt{\pi a_m}.
\]
Subtracting \eqref{eq:actual-upper} from $Q_m$, discarding the nonnegative
tail-area term, and using $P_{<m}\leq|\partial\Omega|$ gives
\begin{align*}
 Q_m-N_{\Omega,R_m}(\theta)
 &\geq4\sqrt\pi\,b_\theta R_m\sqrt{a_m}
 -C_1|\partial\Omega|R_mc_m
 -o(R_m\sqrt{a_m})\\
 &=R_m\sqrt{a_m}\left(4\sqrt\pi\,b_\theta
 -C_1|\partial\Omega|\frac{c_m}{\sqrt{a_m}}-o(1)\right).
\end{align*}
But $c_m/\sqrt{a_m}=\sqrt{a_m}\to0$, while
$R_m\sqrt{a_m}\to\infty$.  Therefore
\begin{equation}\label{eq:unbounded-count-gap}
 Q_m-N_{\Omega,R_m}(\theta)\longrightarrow+\infty.
\end{equation}

\begin{theorem}\label{thm:counterexample}
For the compact regular-closed radial set $\Omega$ in
\eqref{eq:E-Omega}, whose boundary has finite length,
\[
 \limsup_{R\to\infty}R\sup_k\left|\lambda_k^{R\Omega}
 -\frac12\erfc\!\left(
 \sqrt{2\pi}\frac{k-R^2|\Omega|}{R|\partial\Omega|}
 \right)\right|=+\infty.
\]
In particular, Conjecture 2.6 is false as stated.
\end{theorem}

\begin{proof}
If the displayed limsup were finite, the conjectured uniform $O(R^{-1})$
profile would hold.  Lemma \ref{lem:counting}, applied to the fixed set
$\Omega$, would then give
$N_{\Omega,R_m}(\theta)=Q_m+O(1)$.  This contradicts the unbounded gap
\eqref{eq:unbounded-count-gap}.
\end{proof}

\section{Scope, verification, and novelty bounds}

The example exploits infinitely many boundary components at vanishing scales.
It does not contradict the source's Proposition 2.5 for finitely many radial
components.  Nor does it rule out a corrected universality theorem for smooth,
uniformly Ahlfors-regular, or finitely-component boundaries.  The source itself
notes immediately after Conjecture 2.6 that stronger boundary regularity may
be necessary.

The phrase ``have finite boundary'' is read as finite boundary length,
consistent with the quantity $|\partial\Omega|$ in the formula.  This reading
is also distinguished in the source from the explicit ``finite number of
connected components'' hypothesis in Proposition 2.5.  If the conjecture was
intended to assume finitely many boundary components, the present example is a
statement correction rather than a counterexample to that stronger version.

The proof is non-computational.  The checks isolated in the argument are:
\begin{enumerate}[leftmargin=2em]
\item the radial Hermite eigenvalue formula \eqref{eq:gamma-eigs};
\item the Poisson maximal-atom estimate in Lemma \ref{lem:bump};
\item compactness, regular closedness, and the finite perimeter sum
      \eqref{eq:geometry};
\item feasibility of the diagonal induction defining $\ell_m$;
\item the threshold sandwich \eqref{eq:sandwich}; and
\item the profile-to-counting implication in Lemma \ref{lem:counting}.
\end{enumerate}

A bounded novelty search on 9 August 2026 covered the run's registry,
solution, attempt, and proof-gap indexes; arXiv:2502.11805 and the exact paper
title; searches combining ``boundary universality,'' ``localization
operator,'' ``erfc,'' and the displayed scaling; and the author's 2025 thesis.
The thesis still presents the formula as a conjecture.  No later explicit
resolution or infinite-annulus counterexample was found.  This is not an
exhaustive MathSciNet or zbMATH priority search.

\textbf{Human-review recommendation.} High priority.  Verify the uniformity
in the source's finite-annulus Proposition 2.5, the inductive domination of
its prefix constants, and the convention that finite boundary means finite
$\cH^1$ length.  The remaining estimates are elementary and exact.

\sloppy
\begin{thebibliography}{9}
\raggedright
\bibitem{Halvdansson}
S. Halvdansson,
\emph{Empirical plunge profiles of time-frequency localization operators},
arXiv:2502.11805 (2025), Conjecture 2.6 and Proposition 2.5.
\url{https://arxiv.org/abs/2502.11805}.
\end{thebibliography}

\end{document}
